\input{../tex-inputs/latex-leseansicht-vorspann}

\section[Arthur und Olga Schnitzler an Gustav Schwarzkopf, 8. 8. 1911]{L04522 Arthur und Olga Schnitzler an Gustav Schwarzkopf, 8. 8. 1911}
\nopagebreak\mylabel{L04522v}
\rehead{ }\normalsize\beginnumbering\briefempfaengerindex{Schwarzkopf, Gustav@\textsc{Schwarzkopf, Gustav}!zzzSchnitzler, Olga@\emph{von Olga Schnitzler}!1911-08-081@{8. 8. 1911}|(be}\briefempfaengerindex{Schwarzkopf, Gustav@\textsc{Schwarzkopf, Gustav}!zzzSchnitzler, Olga@\emph{von Olga Schnitzler}!1911-08-081@{8. 8. 1911}|(be}\briefempfaengerindex{Schwarzkopf, Gustav@\textsc{Schwarzkopf, Gustav}!zzzSchnitzler, Arthur@\emph{von Arthur Schnitzler}!1911-08-081@{8. 8. 1911}|(be}
\toendnotes[C]{\smallbreak\pagebreak[2]}
\correspDesc{Versand  durch Arthur Schnitzler, Olga Schnitzler,  am 8. 8. 1911 in Semmering
\newline{}Übermittlung  am 8. 8. 1911 in Semmering
\newline{}Erhalt  durch Gustav Schwarzkopf im Zeitraum [9. 8. 1911 – 13. 8. 1911?] in Wien}\toendnotes[C]{\smallbreak}
\Standort{DLA, A:Schnitzler, HS.1985.1.1897.}
\physDesc{Bildpostkarte, 387 Zeichen
\newline{}Handschrift Arthur Schnitzler: Bleistift, deutsche Kurrent
\newline{}Handschrift Olga Schnitzler: Bleistift, lateinische Kurrent
\newline{}Versand: Stempel: »\nobreak{}\oindex{Semmering@\textbf{Semmering}|pwk}Semmering, 8. VIII. 11, XII\nobreak{}«.  }\toendnotes[C]{\smallbreak}
\pstart
           \noindent{}\centering{}{\pb}\textcolor{gray}{\textbf{Semmering\oindex{Semmering@\textbf{Semmering}|pw}. Villa Leibenfrost\oindex{Villa Leibenfrost@\textbf{Villa Leibenfrost}|pw} – Dependancen des Südbahnhotel\oindex{Südbahnhotel [Semmering]@\textbf{Südbahnhotel [Semmering]}|pw} und Kurhaus\oindex{Kurhaus Semmering@\textbf{Kurhaus Semmering}|pw}}}\pend
           \pstart{}{\pb}\textsc{Herrn}\pend{}\pstart{}\textsc{Gustav Schwarzkopf}\pend{}\pstart{}Wien\oindex{Wien@\textbf{Wien}|pw}\pend{}\pstart{}\textsc{I. Tiefer Graben 23}\oindex{Wien@\textbf{Wien}!I., Innere Stadt@\textbf{I., Innere Stadt}!Tiefer Graben 23@\textbf{Tiefer Graben 23}|pw}\pend{}{\bigskip}\vspace{1em}
\pstart
           \noindent{}{\pb}Lieber Guſtav, ſo ſind wir denn hier
      gelandet, nach kurzem \label{K_L04522-1v}\edtext{\textsc{St. Gilgner\oindex{St. Gilgen@\textbf{St. Gilgen}|pw}}
               Aufenthalt}{\lemma{\textnormal{\emph{St. Gilgner
               Aufenthalt}}}\Cendnote{\textnormal{Die Tage vom 24. 7. 1911 bis zum 29. 7. 1911 verbrachten Olga\pwindex{Schnitzler, Olga 17.\,1.\,1882 Wien – 13.\,1.\,1970 Lugano@\textsc{Schnitzler, Olga} (17.\,1.\,1882 Wien – 13.\,1.\,1970 Lugano)|pwk} und Arthur Schnitzler in St. Gilgen\oindex{St. Gilgen@\textbf{St. Gilgen}|pwk}.}}}\label{K_L04522-1}, – hauptſächlich um mit
      Mama\pwindex{Schnitzler, Louise 8.\,7.\,1840 Kőszeg – 9.\,9.\,1911 Wien@\textsc{Schnitzler, Louise} (8.\,7.\,1840 Kőszeg – 9.\,9.\,1911 Wien)|pwv} hier zuſa{\geminationm}en zu ſein, die aber
      gleich wieder nach Baden\oindex{Baden bei Wien@\textbf{Baden bei Wien}|pw}
      gefahren iſt. Abgeſehn von den
      Schrecken der Saison iſt es
      hier ſchön wie i{\geminationm}er u
               {\pb}wir gedenken (mit den Kindern\pwindex{Schnitzler, Heinrich 9.\,8.\,1902 Hinterbrühl – 12.\,7.\,1982 Wien@\textsc{Schnitzler, Heinrich} (9.\,8.\,1902 Hinterbrühl – 12.\,7.\,1982 Wien)|pw}\pwindex{Cappellini, Lili 13.\,9.\,1909 Wien – 26.\,7.\,1928 Venedig@\textsc{Cappellini, Lili} (13.\,9.\,1909 Wien – 26.\,7.\,1928 Venedig)|pw}) auch eine Woche
      zu bleiben.\pend
           \pstart Herzlichſt Ihr \spacefill\mbox{A.}\pend{}
\pstart
           \textsc{Südbahnhotel\oindex{Südbahnhotel [Semmering]@\textbf{Südbahnhotel [Semmering]}|pw}}\hspace*{2.5em}8. 8. 911\pend
           \selectlanguage{ngerman}\vspace{1em}
\pstart
           \noindent{}{[}hs. O. Schnitzler:{]} \spacefill\mbox{Olga.}\pend
           \selectlanguage{ngerman}\endnumbering\briefempfaengerindex{Schwarzkopf, Gustav@\textsc{Schwarzkopf, Gustav}!zzzSchnitzler, Olga@\emph{von Olga Schnitzler}!1911-08-081@{8. 8. 1911}|)be}\briefempfaengerindex{Schwarzkopf, Gustav@\textsc{Schwarzkopf, Gustav}!zzzSchnitzler, Olga@\emph{von Olga Schnitzler}!1911-08-081@{8. 8. 1911}|)be}\briefempfaengerindex{Schwarzkopf, Gustav@\textsc{Schwarzkopf, Gustav}!zzzSchnitzler, Arthur@\emph{von Arthur Schnitzler}!1911-08-081@{8. 8. 1911}|)be}\mylabel{L04522h}
\begin{anhang}
\end{anhang}\newcommand{\dateiname}{L04522}\newcommand{\titel}{Arthur und Olga Schnitzler an Gustav Schwarzkopf, 8. 8. 1911}\newcommand{\editorInnen}{Herausgegeben von Jahnke, SelmaMüller, Martin Anton}\input{../tex-inputs/latex-leseansicht-abspann}
